\documentclass[12pt,a4paper]{article}

\usepackage{fontspec}
\usepackage{geometry}
\usepackage{setspace}
\usepackage{microtype}
\usepackage{parskip}

\setmainfont{Times New Roman}
\geometry{
  a4paper,
  left=4cm,
  right=3cm,
  top=3cm,
  bottom=3cm
}
\pagestyle{empty}
\setlength{\emergencystretch}{2em}

\begin{document}

\begin{center}
{\large\bfseries ABSTRACT}

\vspace{0.45cm}

\textbf{SYSTEM ARCHITECTURE DESIGN FOR AUTOMATED CONTAINER INSPECTION IN PORTS}

\vspace{0.4cm}
by

\vspace{0.4cm}
Aththariq Lisan Quran Daulah Sentono\\
18222013
\end{center}

\begin{spacing}{1.0}
Container inspection processes at ports still face information fragmentation
among stakeholders, inspection records that are not yet connected to container
history, and limited visual evidence that can be referenced later. These
conditions create a risk that container identity, inspection results, and
inspection evidence are scattered across different media. This final project
aims to design a container inspection automation system architecture that
supports container identification, inspection result recording, visual evidence
management, and container history presentation within a single information
structure.

The design was conducted using the Design Science Research approach, which
includes problem identification, objective definition, artifact design,
demonstration through system realization, and result evaluation. The system
architecture was structured using a layered approach consisting of presentation,
application, data, and edge layers. In this design, the edge component handles
video stream acquisition and initial inference, the API service manages
orchestration and data persistence, and the web application provides an
interface for monitoring inspection results. From the data perspective, the
container entity is positioned as the central link between identity, inspection
results, manifest information, and visual evidence.

The system realization shows that the proposed architecture can be implemented
as a flow connecting the edge component, API service, data storage, and web
interface. The evaluation was carried out through integration testing of the
main flows, including the creation of a container entity through OCR, propagation
of the container number to other scanning flows, storage of inspection results
in the core data model, and presentation of results in the web application.
Based on these evaluation results, the proposed design supports container
history recording within the scope of the internal realization of this final
project. The main contribution of this final project is an architecture design
that maps responsibility separation among components, places the container
entity as the center of inspection information, and evaluates the suitability of
the design against the system realization.

Keywords: system architecture, container inspection automation, data
integration, port, container history.
\end{spacing}

\end{document}
